\scanpage{100}
\mspage{- 12 -}\hfill n\textsuperscript{o} 285

% L1
se borner au cas où $I$ est un intervalle de \hspace{1.2em}, et comme la proposition
% L2
est triviale si $I$ est fini, on peut supposer $I=\mathbb{N}$. Soit $x_p\handwritten{\in}G_p$ ;
% L3
l'hypothèse permet de définir par récurrence sur $n\handwritten{\geqslant}p$ une suite $(x_n)_{n\handwritten{\geqslant}p}$
% L4
satisfaisant aux conditions $x_n\handwritten{\in}G_n$, $f_{n,n+1}(x_{n+1})=x_n$ ; il est clair
% L5
alors que si on pose, pour $m\handwritten{<}p$, $x_m=f_{mp}(x_p)$ la suite $(x_n)_{n\handwritten{\geqslant}0}=x$
% L6
appartient à $G$ et que $f_p(x)=x_p$, d'où la proposition.

% L7
PROPOSITION 23 - \uline{Les notations étant celles du cor}.1 \uline{de la prop.}20,
% L8
\uline{on suppose en outre que pour tout} $\alpha\handwritten{\in}I$, \uline{l'application canonique} $\varphi_\alpha$ \uline{de} $A$
% L9
\uline{dans} $A_\alpha$ \uline{soit surjective}. \uline{Alors}, \uline{si pour tout} $\alpha\handwritten{\in}I$ $g_\alpha(A_\alpha)$ \uline{est dense}
% L10
\uline{dans} $B_\alpha$, $g(A)$ \uline{est dense dans} $B$.

% L11
En effet, il suffit de montrer que, pour tout $\alpha\handwritten{\in}I$ et tout ensemble
% L12
ouvert non vide $U_\alpha$ dans $B_\alpha$, il existe $x\handwritten{\in}A$ tel que $\handwritten{\psi}_\alpha(g(x))\handwritten{\in}U_\alpha$ (en
% L13
désignant par $\handwritten{\psi}_\alpha$ l'application canonique de $B$ dans $B_\alpha$). Mais par hypothèse
% L14
il existe $x_\alpha\handwritten{\in}A_\alpha$ tel que $g_\alpha(x_\alpha)\handwritten{\in}U_\alpha$, et il existe $x\handwritten{\in}A$ tel que
% L15
$\varphi_\alpha(x)=x_\alpha$ ; comme $\handwritten{\psi}_\alpha(g(x))=g_\alpha(\varphi_\alpha(x))$, cet élément répond à la question

% L16
PROPOSITION 24 - \uline{Soient} $(G_\alpha,f_{\alpha\handwritten{\beta}})$ \uline{un système projectif de groupes}
% L17
\uline{topologiques}, $G=\underset{\handwritten{\longleftarrow}}{\lim}\,G_\alpha$, $f_\alpha$ \uline{l'application canonique de} $G$ \uline{dans} $G_\alpha$. \uline{On}
% L18
\uline{suppose que} : 1\textsuperscript{o}) \uline{pour tout} $\alpha\handwritten{\in}I$, $f_\alpha$ \uline{est surjective} ; 2\textsuperscript{o}) \uline{pour tout}
% L19
\uline{couple d'indices tel que} $\alpha\handwritten{\leqslant}\handwritten{\beta}$, $f_{\alpha\handwritten{\beta}}$ \uline{est un homomorphisme de} $G_{\handwritten{\beta}}\unsure{{}^{\cdot\cdot}|{:}}$ \uline{sur}
% L20
$G_\alpha$. \uline{Dans ces conditions}, $f_\alpha$ \uline{est un homomorphisme de} $G$ \uline{sur} $G_\alpha$.

% L21
En effet, soit $V$ un voisinage quelconque de $e$ dans $G$. Il existe un
% L22
indice $\handwritten{\beta}\handwritten{\geqslant}\alpha$ et un voisinage $V_{\handwritten{\beta}}$ de $e_{\handwritten{\beta}}$ dans $G_{\handwritten{\beta}}$ tel que $W=f^{-1}_{\handwritten{\beta}}(V_{\handwritten{\beta}})$
% L23
soit un voisinage de $e$ contenu dans $V$. Comme $f_{\handwritten{\beta}}$ est surjective, on a
% L24
$V_{\handwritten{\beta}}=f_{\handwritten{\beta}}(W)$, et comme $f_{\alpha\handwritten{\beta}}$ est un homomorphisme surjectif, $f_{\alpha\handwritten{\beta}}(V_{\handwritten{\beta}})$
% L25
est un voisinage de $e_{\handwritten{\unsure{\delta|\alpha}}}$ dans $G_\alpha$ ; mais on a $f_\alpha(W)=f_{\alpha\handwritten{\beta}}(f_{\handwritten{\beta}}(W))=$
% L26
$=f_{\alpha\handwritten{\beta}}(V_{\handwritten{\beta}})$ d'où la proposition (th.1).

% L27
Soit $N_\alpha$ le noyau de $f_\alpha$ ; les $N_\alpha$ constituent une base de filtre, formée
% archive sheet number: 99
